

\section{Usage Guide}

Various \ac{sihft} passes can be invoked via command line arguments to the LLVM static compiler \texttt{llc} program.
This chapter assumes that the LLVM compiler has been installed as per instructions in earlier sections.
We refer the installation location of this compiler to \texttt{<llvm\_install>} in rest of the chapter.

\subsection{\ac{sihft} Compilation}
Each source file (translation unit) in C/C++ project has to be first compiled to LLVM IR code using \texttt{clang}.
Following command can be used to generate human readable LLVM IR code representation of given C source file.
\begin{framed}
 \begin{lstlisting}[language=bash, basicstyle=\small\ttfamily]
$ <llvm_install>/bin/clang
    -emit-llvm -S <other options>
    main.c -o main.ll
  \end{lstlisting}
\end{framed}

The next step in the compilation is to pass the LLVM IR file to \texttt{llc} program:\\
\begin{lstlisting}[language=bash, basicstyle=\small\ttfamily, frame=single]
$ <llvm_install>/bin/llc
    <SIHFT options> <other options>
    main.ll -o main_hardened.s
\end{lstlisting}

\texttt{<SIHFT options>} refer to various \ac{sihft} transformations that we have implemented. A brief overview of these
options is provided in Table~\ref{tab:sihft-options}.

\begin{table}[htb]
 \centering
 \caption{Supported \ac{sihft} options}
 \label{tab:sihft-options}

 \begin{tabular}{|l|l|}
  \hline
  \textbf{Option} & \textbf{Description}                                     \\
  \hline
  -NZDC=foo,bar   & NZDC transformation on \textit{foo,bar} functions  \\
 
  -NZDC+NEMSIS=foo,bar   & NZDC and NEMESIS transformation \\
  & on \textit{foo,bar} functions  \\
  
  -NZDC+NEMSEC=foo,bar   & NZDC and NEMESEC (NEMESIS w/\\ 
   & store fallbacks) transformation on \\
   &\textit{foo,bar} functions  \\
 
  -SWIFT=foo,bar  & SWIFT transformation 
  on \textit{foo,bar}\\ & functions  \\

  -EDDI=foo,bar & EDDI transformation on \textit{foo,bar} functions \\
  
  -RASM=foo,bar   & RASM transformation 
  on \textit{foo,bar}\\ & functions  \\
  
  -RACFED=foo,bar & RACFED transformation
   on \textit{foo,bar}\\ & functions  \\
  
  -CFCSS=foo,bar  & CFCSS transformation on \textit{foo,bar} functions \\
  
  -FGS=foo,bar    & use fine-grain scheduling on \textit{foo,bar} functions \\
  -REPAIR         & REPAIR transformation on NZDC code                   \\
  \hline
 \end{tabular}
\end{table}

The final step in code generation is to use \texttt{gcc} for RISCV architectures to assemble and link the
assembly code for a particular RISCV processor.

\subsubsection{Stack pointer arguments}
EDDI allocates separate primary and shadow stack frames. For call arguments
derived from the current frame, the argument check compares the shadow address
with the primary address plus the original frame size. The analysis follows
register values across basic blocks and preserves value checks for scalar
arguments loaded from stack memory. SWIFT uses one stack and checks its argument
values directly. Both instruction schedules retain these checks.

Address tracking covers copies and full-width addition and subtraction in integer
registers. Mixed stack and non-stack provenance, or unsupported arithmetic on a
tracked stack address, produces a compilation error at the affected call.
Forwarding pointers from another EDDI frame, spilling pointer values to memory,
and transporting pointer bits through floating-point registers require additional
provenance tracking and remain unsupported by this analysis.

Before a callee establishes its own frame pointer, the incoming \texttt{s0}
and its shadow belong to the caller. EDDI preserves both independently in the
callee-saved spill slots, without requiring their values to be equal: they can
point into separate caller frames whose size is unknown to the callee. The
exception is identified from the saved-register spill metadata. Ordinary data
stores, including values held in \texttt{s0}, retain their value checks.

\subsubsection{Example}
For the sake of example, assume we want to protect the \texttt{crc} program from the MiBench suite. The
compiled program is to run on the RISCV spike simulator. The C project contains the following C files:

\begin{itemize}
 \item{crc.h: header for CRC library functions}
 \item{crc.c: source for CRC library functions}
 \item{main.c: contains the \texttt{main} program to invoke crc library functions}
\end{itemize}

Following commands are issued on bash shell to create the RISCV binary with
\begin{itemize}
 \item NZDC, RASM protection on \texttt{crcSlow} function
 \item RASM protection on \texttt{crcFast} function
 \item NZDC, CFCSS protection on \texttt{main} function
\end{itemize}

\begin{framed}
 \begin{lstlisting}[language=bash, basicstyle=\small\ttfamily]
$ cd <crc-project-source>

$ <clang> -emit-llvm -O2 -S --target=riscv64
    -march=rv64gc -mabi=lp64d
    -isystem <riscv64-gcc>/riscv64-unknown-elf/
    main.c -o main.ll
$ <clang> -emit-llvm -O2 -S --target=riscv64
    -march=rv64gc -mabi=lp64d
    -isystem <riscv64-gcc>/riscv64-unknown-elf/
    crc.c -o crc.ll

$ dmr_funcs="main,crcSlow,crcFast"
$ cfcss_funcs="main"
$ rasm_funcs="crcSlow,crcFast"

$ <llc> -O2 -march=riscv64 -mattr=+m,+a,+d,+c
    -NZDC=${dmr_funcs} -CFCSS=${cfcss_funcs}
    main.ll -o main_hardened.s
$ <llc> -O2 -march=riscv64 -mattr=+m,+a,+d,+c
    -NZDC=${dmr_funcs} -RASM=${rasm_funcs}
    crc.ll -o crc_hardened.s

$ <riscv64-gcc>/bin/riscv64-unknown-elf-gcc -O2
    -march=rv64gc -mabi=lp64d
    main_hardened.s crc_hardened.s -o crc_hardened.elf

$ spike pk crc_hardened.elf

\end{lstlisting}
\end{framed}

% \newpage
% \section{Configuration via TOML script}
